\chapter{Resilience}\label{ch:ll:resilience}

On 1 October 2020, a memory module failed in one of the two devices of a shared disk system inside the Tokyo Stock
Exchange's trading system. Operations ``were supposed to switch automatically'' to the second device, but the automatic
switch did not work for that kind of failure unless a setting had been changed, and it had not. Engineers switched over by
hand at 9:26, less than half an hour after the market should have opened, and every function returned to normal. Trading
nevertheless stayed halted for the whole day. Orders received before the exchange cut its network had been matched, and
their executions had accumulated inside the system without being sent to the participants; only a limited number of
participants could have resent their orders, and there were no agreed, tested rules for restarting. The hardware recovered in
minutes; the state did not. This chapter is about the second problem: how a trading system keeps one agreed history of what
it received and what it sent, so that when a machine fails another can take over in milliseconds without losing an order
or sending one twice.

\section{Failure modes of a trading system}

Components fail in a few ways, and each needs a different answer. A process crashes, or its machine does: it stops, and the
question is who takes over and from what state. A process hangs or slows down: it has not stopped, so it may come back and
act after someone else has taken over, which is worse than a crash. A network link fails: each side may believe the other
is dead while both are alive. A venue session drops: the orders at the venue carry on, may fill or may be cancelled by the
venue (chapter~21), and the firm does not know which until it reconnects. And the \omterm{def:ll:resilience:failover}{failover} machinery itself fails, as in
Tokyo, because it runs only when something else has already gone wrong and is therefore the least exercised code in the
system. The operational risk of One Quant Book~6 (chapter~28) is, for a trading system, mostly this list.

The design goal follows: every input the system acts on, and every order it sends, must be recorded before it matters, in
one order that everyone agrees on, so that any machine can rebuild the state by replaying the record, and any order can be
recognised if it is sent again.

\section{Sequencer architectures}

\begin{definition}[Sequencer architecture]\label{def:ll:resilience:sequencer}
A \emph{sequencer architecture}\index{sequencer architecture} is a design in which every input of a system (market events,
order reports, commands) passes through a single component, the sequencer, that assigns it the next number of a global
sequence and appends it to a durable journal before any other component acts on it; every component then processes the
journal in sequence order.
\end{definition}

The sequencer turns many unordered sources into one ordered history. It is the \omterm{def:ll:logging-capture-and-replay:journal}{input journal} of chapter~23 lifted from one
process to the system: the market data from the \omterm{def:ll:the-feed-handler:handler}{feed handler}, the reports from the \omterm{def:ll:order-gateway-and-order-management:gateway}{order gateway}, the parameter changes from
the control plane all become numbered entries of one journal. The strategy engine's replay of chapter~23 becomes the normal
way of running: a component is a function of the journal, and two components that read the same journal compute the same
thing. The cost is one more hop on the path of every input (a microsecond in the build's simulation, more when the journal
is replicated before it is acknowledged), and a component whose own failure must be handled; the retail exchange whose architecture Fowler described
(chapter~20) kept its business logic on one thread fed by exactly such a journal, from which its state could be rebuilt.

\begin{omfigure}
\begin{tikzpicture}[font=\scriptsize,>=Stealth,box/.style={draw,rounded corners=2pt,minimum height=0.7cm,align=center,
  inner sep=3pt,fill=omDef!12}]
\node[box] (md) at (0,1.0) {market data\\ (feed handler)};
\node[box] (rp) at (0,-1.0) {order reports\\ (gateway)};
\node[box,fill=omThm!18] (sq) at (3.0,0) {sequencer\\ numbers inputs};
\node[box,fill=gray!10,minimum width=2.2cm] (jr) at (6.0,0) {journal\\ epoch $|$ seq $|$ entry};
\node[box] (pr) at (9.2,1.0) {primary replica\\ \textbf{sends}};
\node[box] (bk) at (9.2,-1.0) {backup replica\\ computes, silent};
\node[box,fill=gray!10] (vn) at (12.0,1.0) {venue};
\draw[->] (md) -- (sq);
\draw[->] (rp) -- (sq);
\draw[->,thick] (sq) -- (jr);
\draw[->] (jr) -- (pr);
\draw[->] (jr) -- (bk);
\draw[->,thick] (pr) -- node[above]{orders} (vn);
\draw[->,dashed] (pr) -- node[right]{heartbeats} (bk);
\draw[->,dashed,omThm] (bk.east) -- ++(0.5,0) |- node[pos=0.25,right,align=left]{on failover: new epoch,\\ reconnect, reconcile,\\ resend} (vn.south);
\end{tikzpicture}

\omcaption{A \omterm{def:ll:resilience:sequencer}{sequencer architecture} with a hot backup. Every input is numbered and journaled before anyone acts on it; both
replicas apply the journal and compute the same state and the same orders, but only the primary sends. When the \omterm{def:ll:protocols-i-fix:session}{heartbeats}
stop, the backup raises the journal's epoch, reconnects the order-entry session, reconciles and resends what no report
acknowledges.}\label{fig:ll:resilience:arch}
\end{omfigure}

\section{Replicated state machines and failover}

\begin{definition}[Replicated state machine, primary--backup replication]\label{def:ll:resilience:rsm}
A \emph{replicated state machine}\index{replicated state machine} is a deterministic program run as several copies (replicas)
that apply the same sequence of inputs and therefore pass through the same states and produce the same outputs.
\emph{Primary--backup replication}\index{primary--backup replication} runs one replica as the primary, whose outputs are sent,
and others as backups, whose outputs are computed and discarded, ready to replace the primary.
\end{definition}

\begin{definition}[Failover, split brain]\label{def:ll:resilience:failover}
A \emph{failover}\index{failover} is the replacement of a failed primary by a backup, which takes over its role and its
connections. \emph{Split brain}\index{split brain} is the state in which two replicas both act as primary, typically because a
backup has taken over from a primary that was slow or cut off rather than dead; each sends outputs the other does not know
about.
\end{definition}

The state machine approach is old and general (Schneider's tutorial of 1990), and the hard part of it, agreeing on the next
entry of the log when machines fail, is what consensus algorithms such as Raft (Ongaro and Ousterhout, 2014) solve. A trading
system usually needs less: one sequencer writing the journal, replicated synchronously to a second machine, and a rule for
deciding who is primary. The build's replicas are the strategy state of the book in miniature: \texttt{firm.sequencer}'s
replica trades a hundred shares against every price move with an immediate-or-cancel order, within a position limit, and keeps
its orders by an identifier derived from the journal entry that created them. The Python, C++ and Rust replicas replay the
fixture's journal (a \omterm{def:ll:resilience:failover}{failover}, with a second epoch, resends and a duplicate in it) to the same states, hash included.

\omterm{def:ll:resilience:failover}{Split brain} is prevented by \emph{fencing}. The journal carries an epoch, a number that only grows; each primary writes under
the epoch in which it became primary, and a backup that takes over first raises the journal's epoch. From then on, the journal
refuses any append under an older epoch (\cref{lst:ll:resilience:fence}): a primary that was only slow, and wakes up still
believing itself in charge, can compute what it likes but can no longer record anything, and nothing it has not recorded is
sent. The gateway enforces the same rule on outputs, since the venue does not know about epochs.

\omcode{code/firm/sequencer/cpp/firm_sequencer.hpp}{32}{42}{The journal's fence: after a new epoch, a writer from an older
one cannot append.}\label{lst:ll:resilience:fence}

Failure is detected by \omterm{def:ll:protocols-i-fix:session}{heartbeats}: the primary sends one every millisecond, and the backup declares it dead after three
missed. Detection is therefore the largest part of the build's recovery time: 2.8 to \qty{3.1}{\milli\second} from the failure
to a backup that is primary, reconnected and reconciled, of which the reconnection and the replay of the missed reports take
tens of microseconds. A shorter timeout recovers faster and declares more slow primaries dead, which is exactly when fencing
matters. The backup in the build is hot: it applies every journal entry as it is written, so there is nothing to catch up. A
cold backup must replay the journal first; the C++ replica applies it at about \qty{37}{\nano\second} an entry, so a day's
journal of 37 million entries (chapter~23) takes about a second and a half, acceptable for a restart and not for a \omterm{def:ll:resilience:failover}{failover}.

\section{Exchange disconnects and recovery}

\begin{definition}[Idempotent processing]\label{def:ll:resilience:idempotent}
An operation is processed \emph{idempotently}\index{idempotent processing} if doing it twice has the same effect as doing it
once. For orders, it means that every order carries an identifier fixed when it is created, so that the venue, or the firm's
own gateway, recognises a second copy and discards it.
\end{definition}

The dangerous moment of a \omterm{def:ll:resilience:failover}{failover} is the orders in the gap: those the old primary sent whose reports were never journaled.
The new primary knows they exist, since it computed them from the journal, but not whether they left: the old primary may
have died before sending, while the order was on the wire, or after the venue executed it and before the report came back. It
has three ways to proceed (\cref{lst:ll:resilience:resend}). It can resend them at once under fresh identifiers, the habit of
gateways that never reuse an identifier. It can first log in again, ask the venue for the reports it missed (the sequenced
replay of chapter~21), reconcile, and resend only what is still unaccounted for. Or it can resend at once under the orders'
original identifiers, derived from the journal, and let the venue reject any copy of an order it already has.

\omcode{code/firm/sequencer/firm_sequencer_sim.py}{116}{137}{The takeover: fence the journal, then resend what no report
acknowledges, under the original identifiers or under fresh ones (a resend is journaled, so that the replicas
agree).}\label{lst:ll:resilience:resend}

\Cref{tab:ll:resilience:cuts} kills the primary at every stage of every one of the scenario's 46 orders, against Book 10's
matching engine. With fresh identifiers and no reconciliation, a failure while an order is on the wire or at the venue but not
yet reported sends it twice: 54 of the 184 cut points end with a duplicate, one for each of the 27 orders that filled, in each
of the two stages, and the worst push the position past its limit, to 1\,100 shares against 1\,000. Reconciling first
removes them all in this model, because every order is reported within \qty{40}{\micro\second} and the venue replays its
reports. Idempotent identifiers remove them regardless of the order of operations, which matters when the venue is slow to
report or the session cannot be replayed: the participants of the Tokyo exchange in 2020 were in exactly that position. In
every case the new primary's state equals a fresh replay of the journal and the position the venue holds.

\begin{table}[ht]
\centering\small
\begin{tabular}{lrrrrr}
\toprule
procedure & \multicolumn{4}{c}{duplicates, by stage of the order at the failure} & largest position \\
 & journaled & in flight & at the venue & reported & \\
\midrule
fresh identifiers & 0 & 27 & 27 & 0 & 1\,100 \\
fresh identifiers, reconciled first & 0 & 0 & 0 & 0 & 1\,000 \\
idempotency keys & 0 & 0 & 0 & 0 & 1\,000 \\
\bottomrule
\end{tabular}
\caption{The primary killed at each of four stages of each of the scenario's 46 orders (184 \omterm{def:ll:resilience:failover}{failovers} per procedure) against
the matching engine; recovery took 2.8 to \qty{3.1}{\milli\second} in every case, and the state always equalled a replay of the
journal. Data: \texttt{fig\_failover.py} (deterministic).}\label{tab:ll:resilience:cuts}
\end{table}

\begin{omfigure}
\begin{tikzpicture}[font=\scriptsize,>=Stealth]
\draw[->] (0,0) -- (12.6,0) node[right]{time};
\foreach \x/\l in {0/{input journaled},2.6/{order sent},6.2/{at the venue},9.8/{report journaled}} {
  \draw (\x,-0.12) -- (\x,0.12);
  \node[below,align=center] at (\x,-0.15) {\l};
}
\fill[omDef!15] (0,0.2) rectangle (2.6,1.0);
\node[align=center] at (1.3,0.6) {not sent:\\ a resend is safe};
\fill[omThm!25] (2.6,0.2) rectangle (9.8,1.0);
\node[align=center] at (6.2,0.6) {a failure here: a resend under a fresh identifier\\ duplicates the order if it executed};
\fill[omMeth!15] (9.8,0.2) rectangle (12.4,1.0);
\node[align=center] at (11.1,0.6) {known:\\ nothing to resend};
\draw[<->] (2.6,1.25) -- node[above]{\qty{20}{\micro\second} on the wire, \qty{20}{\micro\second} back} (9.8,1.25);
\end{tikzpicture}

\omcaption{The life of one order and the window in which a \omterm{def:ll:resilience:failover}{failover} can duplicate it: from the moment it leaves to the moment
its report is journaled, two network crossings. With fresh identifiers and no reconciliation, a failure anywhere in the red
window sends the order twice if it had been executed.}\label{fig:ll:resilience:window}
\end{omfigure}

The same reasoning applies to a simple disconnect, without any \omterm{def:ll:resilience:failover}{failover}: the gateway loses its session, logs in again asking
for the reports it has not seen, applies them (the venue's cancel on disconnect may have cancelled resting orders meanwhile),
reconciles against the \omterm{def:ll:order-gateway-and-order-management:dropcopy}{drop copy}, and only then lets the strategy trade. Book 10's simulator provides each piece: sequenced
replay on login, cancel on disconnect, and the rejection of a duplicate client order identifier.

\section{Testing failover}

\omterm{def:ll:resilience:failover}{Failover} code runs when something else has already failed, so it is the least tested code in a system unless it is tested on
purpose, and the Tokyo incident is the standard example of a switch that worked for the failures it was tested against and not
for the one that happened. The build's test is the table: the primary is killed at every stage of every order of a scenario,
and each run is checked against properties that must hold whatever the failure point (the new primary's state equals a replay
of the journal; its position equals the venue's; no order is executed twice when idempotency keys are used). The properties
matter more than the scenario; exhaustive cut points turn a rare event into a routine test. In production, the same idea is
practised by failing over deliberately, on a schedule, so that the procedure is exercised when nothing else is wrong.

\section{Tutorial: kill the primary everywhere}\label{tut:ll:resilience:tut}

\begin{tutorial}
\textbf{Goal.} Run the sequencer, two replicas and a gateway against the matching engine, kill the primary at every stage of
every order, and compare three recovery procedures. \textbf{End state:} \cref{tab:ll:resilience:cuts}, and green replays of
the \omterm{def:ll:resilience:failover}{failover} journal in three languages.
\begin{tutsteps}
\item \textbf{The scenario.} \texttt{firm\_sequencer\_sim.run(None)} runs 60 prices, a millisecond apart, without failure: 46
orders, a position of $-700$, the state equal to a replay.
\item \textbf{Cut points.} \texttt{cut\_points()} lists the middle of each of the four stages of each order; \texttt{python
fig\_failover.py} runs the three procedures at each.
\item \textbf{Replicas in three languages.} The C++ and Rust replicas replay \texttt{data/journal.txt}, a \omterm{def:ll:resilience:failover}{failover} with fresh
identifiers, and reach the states of \texttt{data/expected.txt}; a writer from the old epoch is refused.
\item \textbf{Catch-up.} \texttt{python bench\_sequencer.py} times a replica applying two million entries.
\end{tutsteps}
\textbf{What to change next.} Delay the venue's reports by a millisecond and watch duplicates appear with reconciliation too;
shorten the \omterm{def:ll:protocols-i-fix:session}{heartbeat} timeout and count the \omterm{def:ll:resilience:failover}{failovers} a slow primary would cause.
\end{tutorial}

\section{Build: the sequencer}\label{bld:ll:resilience:sequencer}

\begin{build}
\bfield{Purpose} The backbone that makes the trading path of the book (chapters~18 to 22) restartable and replicable: one
journal of numbered inputs, replicas that compute from it, and a \omterm{def:ll:resilience:failover}{failover} that neither loses nor duplicates an order.
\bfield{Interface} Python reference \texttt{firm\_sequencer}: \texttt{Journal} with \texttt{append} and \texttt{fence};
\texttt{Replica} with \texttt{apply} and \texttt{unacknowledged}; \texttt{cl\_for}; \texttt{to\_line}, \texttt{from\_line};
\texttt{firm\_sequencer\_sim.run} and \texttt{cut\_points} against Book 10's engine. C++20 \texttt{firm::seq}:
\texttt{Journal}, \texttt{Replica}, \texttt{parse}. Rust \texttt{firm\_sequencer}: \texttt{Journal}, \texttt{Replica},
\texttt{parse}.
\bfield{Rules} Nothing acts on an input before it is journaled; replicas are deterministic functions of the journal; epochs
only grow and fence older writers; order identifiers derive from the journal; resends are journaled.
\bfield{Acceptance tests} \texttt{code/firm/sequencer/}: the \omterm{def:ll:resilience:failover}{failover} journal replayed to the same states in three languages;
fencing refuses an old epoch; gaps refused; at every cut point of the scenario and for each procedure, the state equals a replay
and the venue's position; duplicates only with fresh identifiers and without reconciliation.
\bfield{Stretch} A replicated journal with synchronous acknowledgement by the backup; several gateways behind the replicas;
Raft for electing the primary; the gateway of chapter~21 and the \omterm{def:ll:pre-trade-risk-and-kill-switches:gate}{risk gate} of chapter~22 as replicas.
\end{build}

\begin{omsources}
\item Tokyo Stock Exchange, \emph{The Cause of the Recent Failure in arrowhead and Measures Implemented}, 5 October 2020, and
\emph{The Failure of Equity Trading System on October 1, 2020}, report, 19 October 2020.
\item F.~B.~Schneider, ``Implementing fault-tolerant services using the state machine approach: a tutorial'', \emph{ACM
Computing Surveys} 22(4), 1990.
\item D.~Ongaro and J.~Ousterhout, ``In search of an understandable consensus algorithm'', \emph{USENIX Annual Technical
Conference}, 2014.
\end{omsources}

\section{Exercises}

\begin{exercise}[$\star$]\label{exo:ll:resilience:1}
The journal's epoch is 3. The old primary, which became primary in epoch 2, tries to append. What happens, and what happens to
the orders it computed from that entry?
\end{exercise}

\begin{exercise}[$\star$]\label{exo:ll:resilience:2}
\omterm{def:ll:protocols-i-fix:session}{Heartbeats} every millisecond, three missed declare the primary dead, checked at each \omterm{def:ll:protocols-i-fix:session}{heartbeat} tick. What are the shortest
and longest times from a failure to its detection?
\end{exercise}

\begin{exercise}[$\star$]\label{exo:ll:resilience:3}
An order created by journal entry 46 is the first output of that entry. What is its client order identifier in the build, and
why must every replica compute the same one?
\end{exercise}

\begin{exercise}[$\star\star$]\label{exo:ll:resilience:4}
A strategy sends 500 orders a second; each spends \qty{40}{\micro\second} between leaving and having its report journaled, and
90\% of them execute. A failure happens at a random moment. What is the probability that a fresh-identifier resend without
reconciliation duplicates an order?
\end{exercise}

\begin{exercise}[$\star\star$]\label{exo:ll:resilience:5}
A cold backup applies the journal at \qty{37}{\nano\second} an entry. How long does it take to catch up on a day of 37 million
entries, and on the last hour of it? Why keep the backup hot anyway?
\end{exercise}

\begin{exercise}[$\star\star$]\label{exo:ll:resilience:6}
The venue rejects a copy of an order under an identifier it has seen. The backup resends under original identifiers an order
the venue executed. What does the backup receive, and what must its replica do with it?
\end{exercise}

\begin{exercise}[$\star\star\star$]\label{exo:ll:resilience:7}
\emph{Coding.} Delay the venue's reports by a millisecond in \texttt{firm\_sequencer\_sim} and rerun the three procedures.
Which one now duplicates orders, at which stages, and why do idempotency keys still produce none?
\end{exercise}

\begin{exercise}[$\star\star\star$]\label{exo:ll:resilience:8}
\emph{Find the flaw.} ``Our \omterm{def:ll:resilience:failover}{failover} is simple: the backup watches the primary's \omterm{def:ll:protocols-i-fix:session}{heartbeat}, and when it stops, the backup
restarts the strategy from its configuration and sends new quotes.''
\end{exercise}

\section{Problem: The Failover That Doubled the Orders}

\begin{problem}[{Weekend problem --- counting duplicates over every point of failure}]\label{pb:ll:resilience:1}
Use \cref{tab:ll:resilience:cuts} (\texttt{failover.csv}, \texttt{failover\_summary.csv}): 46 orders in 60 milliseconds, four
stages each, three recovery procedures.

\medskip\noindent\textbf{Part I --- The window.}
\begin{enumerate}
\item Name the four stages of an order's life used as cut points, and say for which of them the new primary cannot tell
whether the order left.
\item How long does an order spend in the two dangerous stages in the simulation?
\item Why do only 27 of the 46 orders produce a duplicate in each dangerous stage?
\item Why does a failure while an order is journaled but not yet sent never duplicate it?
\end{enumerate}

\medskip\noindent\textbf{Part II --- Over time.}
\begin{enumerate}[resume]
\item What is the scenario's order rate?
\item If the failure time is uniform over the scenario, what is the probability that a fresh-identifier \omterm{def:ll:resilience:failover}{failover} duplicates
an order?
\item What would it be for a strategy sending 5\,000 orders a second with the same latencies?
\item Why does the largest position reach 1\,100 with fresh identifiers?
\end{enumerate}

\medskip\noindent\textbf{Part III --- The remedies.}
\begin{enumerate}[resume]
\item Why does reconciling first remove every duplicate in this model?
\item Give two real situations in which reconciling first would not be enough.
\item Why do idempotency keys work whatever the order of operations?
\item What must hold for idempotency keys to be possible at all?
\end{enumerate}

\medskip\noindent\textbf{Part IV --- The verdict.}
\begin{enumerate}[resume]
\item State the \emph{named result}: duplicates after a \omterm{def:ll:resilience:failover}{failover} without idempotency keys over the failure points, with them,
and the recovery time.
\item What dominates the recovery time, and what is the price of shortening it?
\item What does fencing add that idempotency keys do not?
\item What went wrong in Tokyo in 2020, in the chapter's terms?
\item Why must the resends themselves be journaled?
\item Which properties does the test check at every cut point, and why properties rather than expected outputs?
\item What would you add to test a cold backup?
\item In one sentence: what makes a \omterm{def:ll:resilience:failover}{failover} safe?
\end{enumerate}
\end{problem}

\section{Interview questions}

\begin{interviewq}[$\star$ \iqroles{developer}]\label{iq:ll:resilience:1}
What is a sequencer, and why would a trading system put one in front of everything?
\end{interviewq}

\begin{interviewq}[$\star\star$ \iqroles{developer}]\label{iq:ll:resilience:2}
Your primary stops sending \omterm{def:ll:protocols-i-fix:session}{heartbeats}. How does the backup take over without sending duplicate orders?
\end{interviewq}

\begin{interviewq}[$\star\star$ \iqroles{developer}]\label{iq:ll:resilience:3}
What is \omterm{def:ll:resilience:failover}{split brain}, and how do you prevent it?
\end{interviewq}

\begin{interviewq}[$\star\star$ \iqroles{developer, trader}]\label{iq:ll:resilience:4}
Your session to the exchange drops for five seconds. What do you do, in order, when it comes back?
\end{interviewq}

\begin{interviewq}[$\star\star$ \iqroles{developer}]\label{iq:ll:resilience:5}
How would you test that \omterm{def:ll:resilience:failover}{failover} works, given that it is needed rarely?
\end{interviewq}

\begin{interviewq}[$\star\star\star$ \iqroles{developer}]\label{iq:ll:resilience:6}
Design a trading system that survives the loss of any one machine with at most a few milliseconds of interruption and no
lost or duplicated order. What are the components, the protocol between them, and the tests?
\end{interviewq}
